\section*{The Propers: Themes and Movement}
\label{triptych:brief-synthesis:start}
\begingroup
\setlength{\parskip}{0.32em}

Continuing mercy gives life and sustains a community capable of charity.
The Church who asks to be cleansed and fortified receives Christ's gift
within a life of gentle correction, persevering service, thanksgiving, and
hope for bodily fulfillment.

\subsection*{Mercy takes the initiative\enspace\properrefs{Int., Coll., Gosp., Off.}}

The Introit directs need toward God through commands to hear, save, and show
mercy. Its separate psalm verse joins the lifting of the soul to a request
for joy. The complete Psalm~85 places that prayer within the Lord's
compassion and the nations' worship, yet ends amid continuing danger.
Augustine hears Christ praying with his members: the Church's need belongs
to the Head who makes their voice his own, and holiness itself is received
from God (\emph{Enarrationes} 85.1--6, Tweed's historical English). Petition
rests upon divine goodness rather than an achievement that merits payment.

The Collect makes the same dependence corporate. Continuing mercy must
cleanse and fortify the Church, who cannot stand safely without God and must
always be governed by his gift. In the Gospel, mercy acts before any request
is narrated. Jesus sees the widow, is moved with compassion, touches the
bier, and commands her son to rise (Luke~7:11--16). The mother receives a
living son. Her particular bereavement remains at the heart of the miracle,
even when the Fathers open its ecclesial meaning.

The Offertory answers the Introit's appeal for hearing with prayer heard and
a new canticle received. Its complete psalm returns from thanksgiving to
need; deliverance has not made further petition unnecessary. Augustine's
Psalm~39 exposition likewise places rescue beside continued walking and a
future crown (39.2--5). Dependence belongs to life after restoration as well
as before it. Mercy establishes the faithful on a way they must still walk.

\subsection*{Restored persons become bearers of charity\enspace\properrefs{Coll., Ep., Gosp., Comm.}}

Paul begins with life in the Spirit and requires conduct corresponding to
it. Vainglory, provocation and envy deform community; gentle correction
seeks the fallen person's restoration while remembering the corrector's
own exposure to temptation. Chrysostom calls that mildness a spiritual
gift and explains reciprocal endurance without demanding every gift of
every person (\emph{Galatians}, chs.~5--6, Alexander's English). Christian
freedom has already taken the form of service through love in Galatians~5.
The final chapter brings that service into the difficult encounter with
another's fault.

Challoner's command, ``Bear ye one another's burdens'' (Gal.~6:2), stands
beside each person's own reckoning. Augustine's Latin \emph{Frangipane V}
§§2--4 opening distinguishes shared weakness from personal answerability;
one's own account includes neglect of a brother's wound (Morin,
pp.~213--215). Aquinas develops endurance of defects, practical relief, and
intercessory satisfaction (\emph{Super Galatas} 6.1). Mercy thus requires
real assistance and responsible self-examination together.

Theodoret gives correction its reason in the corrector's own mutability:
someone exposed to change and temptation must tend the fallen with care.
He identifies Christ's law with the command of charity and closes with
the heavenly Father's beneficence (Gal.~6:1--10, PG~82:499A--502B).
The same passage moves from correction to support of teachers and from
sowing to good for all. Receiving instruction therefore has material
consequences, just as receiving a restored brother has communal ones.
Chrysostom sees support of teachers cultivating humility on both sides
(ch.~6). Their different gifts create reciprocal obligations. The Church
preserved by the Collect's mercy is such a community: persons receive,
contribute, examine their conduct, and take responsibility for one another.

Christ returns the son to his mother. Augustine receives that bodily
raising as the restoration of public sinners to Mother Church
(\emph{Sermon}~98.4--5, MacMullen's English); Bede's continuous Latin
transcription makes return to ecclesial communion explicit
(\emph{In Lucam} II.7, PL~92:418D). Luke names an actual mother and son;
the Fathers' ecclesial reading describes the community into which penitents
are restored. \emph{The Liturgical Year}'s 1909 continuation joins this
Sunday's burden-bearing to Eucharistic mutual abiding, citing the
contextual John~6:57 (p.~348). The appointed Communion is the earlier final
clause of John~6:52: Christ's flesh given for the world's life.

Ambrose's continuous Latin transcription gives the mother's tears a place
within this restorative action. Mother Church intercedes for each sinner
as for an only son, and the touched wood signifies salvation through the
cross (\emph{In Lucam} V.89--92). Bede develops the awakened conscience,
compunction and visible renewed life that end in restored communion
(II.7, PL~92:418A--D). Mercy thus reaches both the wounded person and the
body that receives him.

\clearpage
\label{triptych:brief-synthesis:end}

\subsection*{Life is sustained in conduct and hope\enspace\properrefs{Ep., Grad., Gosp., Sec., Comm., Postcomm.}}

The Gospel restores bodily life; the Communion names the flesh of Christ
as a gift for life; the Epistle directs sowing toward an eternal harvest.
Cyril gives the Gospel--Communion relation a precise theological argument.
His \emph{John} IV.2 explains the appointed gift through flesh united to
the life-giving Word (Pusey's English, pp.~409--414). At pp.~418--419,
commenting on a following verse, he invokes Naim: bodily touch gives life,
and sacramental participation communicates that same life-giving flesh.
Present restoration and the promised blessed life have different extents;
Cyril distinguishes the resurrection of all from the beatitude of the holy.

The body therefore belongs within salvation. Galatians' sowing in one's own
flesh describes disordered conduct within the letter's vice-and-virtue
contrast. John's gift is Christ's own flesh. The Postcommunion expressly
asks the heavenly gift's operation to possess bodies as well as minds.
Schuster reads its efficacy as the faculties brought under grace
(\emph{Sacramentary} III, p.~141). Divine action precedes and directs
conduct without abolishing the Epistle's work of restoring, examining,
sharing, persevering, and doing good.

That conduct remains exposed to temptation. The Secret asks the sacraments
to guard against demonic assaults. Aquinas explains Eucharistic protection
through inward union with Christ and the outward defense of his Passion,
while retaining the recipient's capacity to sin (\emph{Summa} III.79.6).
Augustine distinguishes outward reception from living fellowship and abiding
in Christ (\emph{John} 26.15--18, Gibb's English). The gift nourishes charity;
knowingly remaining in mortal sin obstructs its saving effect
(Aquinas, III.79.3).

Chrysostom makes the moral demand follow his exposition of sacramental
incorporation. His \emph{John} Homily~46 moves from the offered flesh and
manna's fulfillment to restraint of wrath and speech and to works that
accompany faith (§§2--4, Marriott's English). The grace received reaches
the same speech and dealings with others that Paul orders away from
provocation and envy. Aquinas similarly distinguishes the bodily conduct
that receives the gift's effect now from incorruption and glory in the
future (III.79.1, ad~3). The Postcommunion's minds and bodies name agents
whose actions can already be governed by grace while their bodily
fulfillment remains ahead.

Paul's future harvest requires good done while opportunity lasts.
Psalm~91's opening praise belongs to a song in which righteous fruitfulness
outlasts the brief flourishing of the wicked. Augustine receives morning
and night as prosperity and adversity (91.1--4); Theodoret reads continual
praise of beneficence and just judgment (PG~80:1615C--1618A). Praise can
persist through conditions that change. Aquinas similarly distinguishes
capacity for glory from immediate entrance into it (III.79.2).

\subsection*{Particular mercy opens into public praise\enspace\properrefs{Int., Grad., All., Gosp., Off., Ep., Comm.}}

At Naim the son speaks, while the crowd glorifies God and acclaims a great
prophet and divine visitation. Luke records no words of the son. The
crowd's visitation language resonates with the Benedictus and Israel's
worship when God sees affliction (Luke~1:68--79; Exod.~4:29--31).
The Offertory's received canticle makes thanksgiving an act the assembly
can itself perform. \emph{The Liturgical Year} explicitly joins Mother
Church's restored children to this chant of thanksgiving (1909, p.~354).

The Alleluia places such particular mercy within universal kingship.
Its \latin{super omnem terram} differs from the Clementine Psalm~94:3's
\latin{super omnes deos}. Theodoret's Psalm~94 lemma has the same all-earth
form as the chant (PG~80:1641A--1642A); Psalm~46:3 also carries the phrase.
Augustine receives that Psalm~46 clause as the kingship of the crucified
and ascended Christ, embracing believing Israel and the nations
(\emph{Enarrationes} 46.2--3, Tweed's English).

Such praise can accompany both deliverance and an affliction still borne.
The Gradual names mercy and truth across morning and night; the complete
Introit and Offertory psalms continue to petition even after confessing
God's goodness. Bellarmine receives Augustine's prosperity-and-adversity
reading of Psalm~91 alongside the difference between mercy visible in
creation and judgment praised in faith's darkness (O'Sullivan,
pp.~289--290). Thanksgiving therefore does not require that every trouble
has ended. In the liturgical movement, the crowd's praise for one bodily
restoration joins the Church's song as she seeks continuing protection
and awaits the eternal harvest named by Paul.

The world's breadth does not dissolve the person. One servant prays, one
mother receives her son, and each Christian answers for his conduct.
Paul asks good for all with particular care for the household of faith;
Aquinas grounds universal beneficence in the common divine image while
preserving practical priority (\emph{Super Galatas} 6.2). Mercy reaches
persons within communion and sends them toward others. Christ's gift for
the world's life sustains that charity and directs it toward its promised
fulfillment.

\endgroup
\clearpage
